% Part: first-order-logic
% Chapter: models-theories
% Section: expressing-relations

\documentclass[../../../include/open-logic-section]{subfiles}

\begin{document}

\olfileid{fol}{mat}{exr}

\olsection{\article{structure} \printtoken{S}{structure} وچ تعلقاں دا اظہار}

\begin{explain}
!!{formula}s دا اک وڈا ورتارا ایہ اے کہ !!a{structure}~$\Struct M$ دیاں خاصیتاں
تے تعلقاں نوں زبان~$\Lang L$ دیاں بنیادی علامتاں راہیں ظاہر
کیتا جا سکدا اے، جیہڑی~$\Struct M$ دی زبان اے۔ ایہدے نال ساڈی مراد ایہ اے:
$\Struct M$ دا !!{domain} اشیا دا سیٹ اے۔ !!{constant}s،
!!{function}s تے !!{predicate}s دی~$\Struct M$ وچ تشریح بالترتیب
$\Domain M$ دے کجھ عناصر، $\Domain M$ اُتے فنکشناں تے
$\Domain M$ اُتے تعلقاں راہیں ہوندی اے۔ مثال لئی، جے
$\Obj A^2_0$، $\Lang L$ وچ ہووے، تاں $\Struct M$ ایہدے نال
تعلق~$R = \Assign{{\Obj A^2_0}}{M}$ جوڑدی اے۔ فیر فارمولا
$\Atom{\Obj A^2_0}{\Obj v_1, \Obj v_2}$ ہی اس تعلق دا
\emph{اظہار کردا اے}، ایس معنی وچ: جے متغیراں دی قدر بندی~$s$،
$\Obj v_1$ نوں $a \in \Domain{M}$ تے
$\Obj v_2$ نوں $b \in \Domain M$ ول بھیجے، تاں
\[
Rab \text{\quad جے تے صرف جے\quad} \Sat{M}{\Atom{\Obj A^2_0}{\Obj v_1, \Obj v_2}}[s].
\]
ایتھے متغیراں دی قدر بندی لانی ضروری اے: اسیں ایہ نہیں کہہ
سکدے کہ ”$Rab$ جے تے صرف جے
$\Sat{M}{\Atom{\Obj A^2_0}{a, b}}$“، کیونکہ $a$ تے
  $b$ ساڈی زبان دیاں علامتاں نہیں، بلکہ
  $\Domain{M}$ دے !!{element}s نیں۔

ساڈے کول صرف ایٹمی !!{formula}s ای نہیں؛ اسیں انہاں نوں منطقی
رابطیاں تے مسوّراں نال جوڑ وی سکدے آں۔ ایس لئی ودھ پیچیدہ
!!{formula}s اوہ ہور تعلق تعریف کر سکدے نیں جیہڑے
$\Struct M$ وچ سدھے بنیادی تعلقاں دے طور اُتے شامل نہیں۔
ساڈی دلچسپی ایس وچ اے کہ ایہ کم کیویں ہوندا اے، تے خاص طور اُتے
ایہ کہ کسے !!a{structure} وچ کیہڑے تعلق تعریف کیتے جا سکدے نیں۔
\end{explain}

\begin{defn}
فرض کرو $!A(\Obj v_1,\dots, \Obj v_n)$، $\Lang L$ دا
!!a{formula} ہووے جس وچ صرف $\Obj v_1$،\dots،
$\Obj v_n$ آزاد آون، تے $\Struct M$، $\Lang L$ لئی
!!a{structure} ہووے۔ $!A(\Obj v_1,\dots, \Obj v_n)$
تعلق~$R \subseteq \Domain M^n$ دا \emph{اظہار کردا اے} جے تے
صرف جے
\[
Ra_1\dots a_n \text{\quad جے تے صرف جے\quad} \Sat{M}{\Atom{!A}{\Obj
    v_1,\dots, \Obj v_n}}[s]
\]
ہر اوہ متغیر قدر بندی~$s$ لئی، جیہدے وچ
$s(\Obj v_i) = a_i$ ($i = 1, \dots, n$)، ایہ برابری قائم ہووے۔
\end{defn}

\begin{ex}
حساب دے معیاری نمونے~$\Struct N$ وچ !!{formula}
$\Obj v_1 < \Obj v_2 \lor \eq[\Obj v_1][\Obj v_2]$،
تعلق $\le$ دا اظہار کردا اے، جیہڑا $\Nat$ اُتے اے۔ !!{formula}
$\eq[\Obj v_2][\Obj v_1']$ جانشین تعلق دا اظہار کردا اے؛
یعنی اوہ تعلق $R \subseteq \Nat^2$ جتھے $Rnm$ ایس صورت وچ
قائم اے جدوں $m$، $n$ دا جانشین ہووے۔ فارمولا
$\eq[\Obj v_1][\Obj v_2']$ پیشرو تعلق دا اظہار کردا اے۔
!!{formula}s
$\lexists[\Obj v_3][(\eq/[\Obj v_3][\Obj 0] \land
  \eq[\Obj v_2][(\Obj v_1 + \Obj v_3)])]$ تے $\lexists[\Obj
  v_3][\eq[(\Obj v_1 + {\Obj v_3}')][v_2]]$ دونے
$\Obj <$ تعلق دا اظہار کردے نیں۔ ایس دا مطلب اے کہ حساب دی زبان
وچ محمول علامت~$<$ اصل وچ لازمی نہیں؛ ایہنوں تعریف کیتا جا سکدا اے۔
\end{ex}

\begin{explain}
ایہ خیال صرف خاص !!{structure}s وچ ای دلچسپ نہیں، بلکہ عام طور اُتے
وی اوہدوں اہم اے جدوں اسیں کسے زبان نال اک مطلوب نمونے یا نمونیاں
دا بیان کردے آں، یعنی جدوں اسیں نظرییاں اُتے غور کردے آں۔
ایہ نظریے اکثر صرف تھوڑے جیہے !!{predicate}s نوں بنیادی
علامتاں دے طور اُتے رکھدے نیں، پر جیہڑے دائرے دا اوہ بیان کردے
نیں، اوہدے وچ ہور بہت سارے تعلق وی اہم ہوندے نیں۔ جے ایہ ہور
تعلق، زبان دے بنیادی !!{predicate}s دی تشریح کرن والے تعلقاں
راہیں باقاعدہ طور اُتے ظاہر کیتے جا سکن، تاں اسیں کہندے آں کہ
انہاں نوں اس زبان وچ \emph{تعریف} کیتا جا سکدا اے۔
\end{explain}

\begin{prob}
$\Lang L_A$ وچ اوہ !!{formula}s لبھو جیہڑے ہیٹھلے تعلقاں
دی تعریف کردے نیں:
\begin{enumerate}
\item $n$، $i$ تے $j$ دے درمیان اے؛
\item $n$، $m$ نوں پورا ونڈدا اے (یعنی $m$، $n$ دا مضرب اے)؛
% PNB-MAT-002: اصل دی وضاحت اک نوں وی شامل کردی سی؛ اولی عدد اک توں وڈا ہونا ضروری اے۔
\item $n$ اولی عدد اے (یعنی اوہ اک توں وڈا اے تے $1$ تے $n$ توں علاوہ ہور کوئی عدد
  $n$ نوں پورا نہیں ونڈدا)۔
\end{enumerate}
\end{prob}

\begin{prob}
فرض کرو فارمولا $!A(\Obj v_1, \Obj v_2)$، تعلق
$R \subseteq \Domain M^2$ دا اظہار !!a{structure}~$\Struct M$
وچ کردا اے۔ اوہ
فارمولے لبھو جیہڑے ہیٹھلے تعلقاں دا اظہار کردے نیں:
\begin{enumerate}
\item اُلٹا تعلق $R^{-1}$، جیہڑا $R$ دا اُلٹا اے؛
\item نسبتی حاصل ضرب $R \mid R$؛
\end{enumerate}
کی تسی $R^+$، یعنی $R$ دی تعدی تکمیل، دا اظہار کرن دا کوئی
طریقہ لبھ سکدے او؟
\end{prob}

\begin{prob}
فرض کرو $\Lang{L}$ اوہ زبان اے جیہدے وچ صرف دو تھانواں
والی محمول علامت $<$ ہووے (ہور کوئی !!{constant}s،
!!{function}s یا !!{predicate}s نہ ہون---البتہ
$\eq$ ضرور ہووے)۔ $\Struct{N}$ اوہ ساخت ہووے جیہدے لئی
$\Domain{N} = \Nat$ تے
$\Assign{<}{N} = \Setabs{\tuple{n,m}}{n <
  m}$۔ ہیٹھلے دعویاں دا ثبوت دیو:
\begin{enumerate}
\item $\{ 0 \}$، $\Struct{N}$ وچ تعریف دے قابل اے؛
\item $\{ 1 \}$، $\Struct{N}$ وچ تعریف دے قابل اے؛
\item $\{ 2 \}$، $\Struct{N}$ وچ تعریف دے قابل اے؛
\item ہر $n \in \Nat$ لئی سیٹ $\{ n \}$،
  $\Struct{N}$ وچ تعریف دے قابل اے؛
\item $\Domain{N}$ دا ہر محدود ذیلی سیٹ
  $\Struct{N}$ وچ تعریف دے قابل اے؛
\item $\Domain{N}$ دا ہر ہم متناہی ذیلی سیٹ
  $\Struct{N}$ وچ تعریف دے قابل اے (جتھے
  $X \subseteq \Nat$ ہم متناہی اے جے تے صرف جے
  $\Nat \setminus X$ محدود ہووے)۔
\end{enumerate}
\end{prob}


\end{document}
